\documentclass[11pt,a4paper]{article}

\usepackage[margin=2.5cm]{geometry}
\usepackage{amsmath}
\usepackage{amssymb}
\usepackage{booktabs}
\usepackage{longtable}
\usepackage{enumitem}
\usepackage[hidelinks]{hyperref}
\usepackage{caption}

\newcommand{\ptl}{\ensuremath{\mu\mathrm{g\,L^{-1}}}}
\newcommand{\um}{\ensuremath{\mu\mathrm{m}}}

\title{Reproducibility of the Modified Gutzeit Method for\\
Inorganic Arsenic in Drinking Water\\
\large A quantitative review of reagent stoichiometry, reaction kinetics,\\
gas-handling geometry, and reflectance photometry}
\author{Prepared as a technical working document}
\date{}

\begin{document}
\maketitle

\begin{abstract}
\noindent
This document reviews the chemistry of a two-reagent modified Gutzeit field
test for inorganic arsenic (L-tartaric acid with trace potassium permanganate;
zinc dust with a trace transition-metal rate enhancer; mercuric bromide
indicator pad), and analyses the sources of the large run-to-run variation
observed with duplicate samples. The principal conclusions are:
(i)~neither reagent is stoichiometrically limiting with respect to arsenic,
by roughly five orders of magnitude;
(ii)~the \emph{buffer capacity} of the tartaric acid, not its total molarity,
sets the duration of the productive pH window, and permits only
$\sim\!0.15\,$g of zinc to dissolve before the pH leaves the range in which
arsine generation is rapid;
(iii)~the measurement is read while the reaction is still climbing, so every
rate-determining variable maps directly onto the reported concentration;
(iv)~the delivery of arsine to the indicator pad, not the chemistry, is the
dominant variance source in the as-supplied geometry, because capture
efficiency is flow-dependent;
(v)~a fixed, shim-defined gap over a back-sealed pad is analytically superior
to a force-balanced flapper valve, giving near-quantitative capture with an
order of magnitude more linear range and no force calibration; and
(vi)~multispectral reflectance analysis, correctly transformed through
Kubelka--Munk, can separate arsenic loading from pad wetting and from local
indicator saturation, but only if measured per-pixel rather than
area-averaged.
Derivations, assumptions, and a revised protocol are given.
\end{abstract}

\tableofcontents

\newpage

%======================================================================
\section{Scope and system description}
%======================================================================

\subsection{The kit as supplied}

\begin{description}[leftmargin=2.2cm,style=nextline]
  \item[Reagent 1] L-tartaric acid (CAS 87-69-4), 98.9\%, with
        potassium permanganate $<1$\%.
  \item[Reagent 2] Zinc powder (CAS 7440-66-6), $>99$\%, with other
        metals (proprietary) $<1$\%.
  \item[Indicator] Paper pad, $5 \times 6.5$~mm, bearing $\approx 1$~mg
        mercuric bromide (HgBr$_2$), bonded to a plastic carrier strip.
\end{description}

\noindent
Manufacturer's procedure: add Reagent~1 to 100~mL of sample in a plastic
bottle; mix 15~s; stand 2~min; add Reagent~2; shake 5~s; place the pad in
the headspace and cap; stand 7~min; read. Headspace 97~mL.

\subsection{Reagent provenance}

The formulation corresponds closely to the reagent system claimed in
US~11{,}346{,}825~B2 (Industrial Test Systems, \emph{Arsenic analysis}),
which claims a first particulate reagent comprising a di- or tri-carboxylic
acid with a C--C moiety between the carboxyl groups (tartaric, malic,
succinic, citric) combined with $0.3$--$5$~wt.\% alkali metal permanganate,
and a second particulate reagent comprising zinc with a transition-metal
additive (Fe, Co, Ni, Cu, Sn, W, Mo) that increases the rate of arsine
production. Worked Examples 10, 11 and 14 of that patent use
$2.7$--$2.8$~g of a KMnO$_4$/tartaric acid blend per 100~mL, matching the
sachet mass measured here.

The rate-enhancer concept derives from US~6{,}696{,}300, which claims
acceleration of the detection chemistry by iron and nickel salts.
US~6{,}576{,}143 gives the balanced arsine-generation equations and the
description of the mercuric bromide stain as ``mixed arsenic mercury
halogenides'' producing a yellow-to-tan-to-brown colour change.

The undisclosed $<1$\% metal component of Reagent~2 is therefore almost
certainly a transition-metal salt (sodium tungstate, potassium molybdate,
or an Fe/Ni salt). This matters: \textbf{the catalyst dose is a fixed
percentage of the sachet mass, so it varies exactly as the zinc mass does.}

\subsection{Modifications already adopted by the operator}

Sample mass recorded to 0.01~g; temperature recorded to 1~$^\circ$C;
125~mL Erlenmeyer flask on a magnetic stirrer (73~mL headspace); stir 25~s,
stand 1:35; add Reagent~2, close immediately, stir 1~min, wait 6~min, read.
Analysis is by a purpose-built reflectance colorimeter using a digital
camera with per-Bayer-layer readout under 4--8 narrow-band LEDs.

%======================================================================
\section{Chemistry}
%======================================================================

\subsection{Step 1: acidification and removal of oxidisable interferents}

L-tartaric acid, $\mathrm{C_4H_6O_6}$, $M = 150.087$~g\,mol$^{-1}$, is a
diprotic weak acid with $\mathrm{p}K_{a1} = 2.98$ and
$\mathrm{p}K_{a2} = 4.34$ at $25\,^\circ$C. Dissolving it establishes the
acidic environment for hydride generation. The patent specifies an optimum
reaction pH of approximately $1.8$--$2.5$, and notes that a pH significantly
above $2.5$ reduces the rate of arsine production.

The trace permanganate oxidises dissolved sulfide to sulfate:
\begin{equation}
  5\,\mathrm{H_2S} + 8\,\mathrm{MnO_4^-} + 14\,\mathrm{H^+}
  \longrightarrow 5\,\mathrm{SO_4^{2-}} + 8\,\mathrm{Mn^{2+}}
  + 12\,\mathrm{H_2O}
\end{equation}
This step is essential because hydrogen sulfide reacts with mercury to give
mercury sulfide, which both destroys the pad's sensitivity to arsenic and
produces a grey cast that obscures the arsenic colour. The patent states
that $0.01$--$0.08$~g of an alkali metal permanganate suffices to remove up
to about $4$--$6$~ppm of hydrogen sulfide, consistent with the
$\leq 0.029$~g present here (Section~\ref{sec:inventory}).

Permanganate is also consumed by tartaric acid over the standing period.
Any permanganate surviving into Step~2 will oxidise As(III) to As(V), and
As(V) is reduced to arsine more slowly than As(III) under field-kit
conditions. Variability in the surviving permanganate is therefore a
speciation-mediated rate variable, and is one reason the 2-minute standing
period should not be shortened.

\subsection{Step 2: arsine generation}

Zinc dissolves in the acid, generating hydrogen:
\begin{equation}
  \mathrm{Zn} + 2\,\mathrm{H^+} \longrightarrow
  \mathrm{Zn^{2+}} + \mathrm{H_2}
  \label{eq:zn}
\end{equation}
and the nascent hydrogen reduces dissolved arsenic to arsine:
\begin{align}
  \mathrm{H_3AsO_4} + 4\,\mathrm{Zn} + 8\,\mathrm{H^+}
    &\longrightarrow \mathrm{AsH_3} + 4\,\mathrm{Zn^{2+}}
    + 4\,\mathrm{H_2O} \\
  \mathrm{H_3AsO_3} + 3\,\mathrm{Zn} + 6\,\mathrm{H^+}
    &\longrightarrow \mathrm{AsH_3} + 3\,\mathrm{Zn^{2+}}
    + 3\,\mathrm{H_2O}
\end{align}
The excess hydrogen from~\eqref{eq:zn} is not waste: it is the carrier gas
that strips arsine from solution and transports it to the indicator.
Tartrate additionally complexes Zn$^{2+}$, which helps keep the metal
surface free of passivating hydroxide as the pH rises.

Two well-documented kinetic facts govern this step. First, zinc particle
size controls the kinetics of hydrogen production, and generation must
continue long enough to expel all the arsenic from the sample. Second,
As(III) forms the hydride more rapidly than As(V), so incomplete or variable
pre-oxidation shifts the effective rate.

\subsection{Step 3: stain formation}

Arsine displaces bromide from mercuric bromide in a stepwise sequence.
Mercuric bromide is colourless; the mono-substituted product is yellow, and
the colour becomes progressively browner as further bromide is replaced:
\begin{align}
  \mathrm{AsH_3} + \mathrm{HgBr_2} &\rightarrow
    \mathrm{AsH_2(HgBr)} + \mathrm{HBr}
    &&\text{(yellow)}\\
  \mathrm{AsH_3} + 2\,\mathrm{HgBr_2} &\rightarrow
    \mathrm{AsH(HgBr)_2} + 2\,\mathrm{HBr}
    &&\text{(yellow)}\\
  \mathrm{AsH_3} + 3\,\mathrm{HgBr_2} &\rightarrow
    \mathrm{As(HgBr)_3} + 3\,\mathrm{HBr}
    &&\text{(brownish-yellow to brown)}
  \label{eq:tri}
\end{align}
Rahman \emph{et al.} assign the disubstituted species as yellow and the
trisubstituted as brownish-yellow. At extreme loading, mercury arsenide
(Hg$_3$As$_2$) and metallic deposits can form; these are spectrally neutral
(grey-black) rather than yellow, which has consequences for multispectral
interpretation (Section~\ref{sec:optics}).

The limiting stoichiometry for capacity calculations
is~\eqref{eq:tri}: three HgBr$_2$ per arsenic.

\subsection{Interferences}

\begin{itemize}[nosep]
  \item \textbf{Sulfide} --- the principal interferent; removed by
        permanganate. Note that the standard mercuric bromide paper
        documentation warns that hydrogen sulfide can itself produce a
        yellow colour, so sulfide interference is not reliably
        distinguishable by hue alone.
  \item \textbf{Antimony and selenium} --- also form hydrides, but are
        rarely present in groundwater at interfering levels.
  \item \textbf{Organic arsenic} --- generally not detected without a
        preliminary oxidation step; kit results will read low relative to
        ICP-MS if organoarsenicals are present.
  \item \textbf{Photobleaching} --- the yellow product fades, especially in
        sunlight, so the pad must be read promptly and consistently.
\end{itemize}

%======================================================================
\section{Quantitative reagent inventory}
\label{sec:inventory}
%======================================================================

\subsection{Sachet mass statistics}

Thirty-two sachets of each reagent were weighed net of packaging.

\begin{table}[h]
\centering
\caption{Sachet mass statistics ($n=32$ each).}
\begin{tabular}{lrr}
\toprule
Statistic & Reagent 1 (acid) & Reagent 2 (zinc) \\
\midrule
Mean (g)                    & 2.896 & 1.266 \\
Median (g)                  & 2.887 & 1.298 \\
Standard deviation (g)      & 0.111 & 0.252 \\
Coefficient of variation    & 3.8\% & \textbf{19.9\%} \\
SD excluding high outlier   & 0.057 & --- \\
CV excluding high outlier   & 2.0\% & --- \\
Minimum (g)                 & 2.748 & 0.553 \\
Maximum (g)                 & 3.448 & 1.678 \\
Max/min ratio               & 1.25  & \textbf{3.03} \\
Interquartile range (g)     & ---   & 1.116--1.449 \\
\bottomrule
\end{tabular}
\end{table}

Reagent~1 is well controlled apart from a single sachet 19\% above the mean.
Reagent~2 exhibits a 20\% coefficient of variation and a threefold range,
which is a filling-process defect rather than ordinary dispensing variation.

\emph{Note on the data:} the two columns were recorded in sorted order, so
the pairing between the acid and zinc sachet used in any given test is lost.
Future weighings should record matched pairs so that colour can be regressed
against both masses independently.

\subsection{Molar inventory per test}

Taking the outlier-excluded acid mean of 2.878~g:
\begin{align}
  n(\mathrm{H_2T}) &= \frac{2.878 \times 0.989}{150.087}
     = 1.896\times10^{-2}\ \mathrm{mol}
     \quad\Rightarrow\quad C_T = 0.190\ \mathrm{M} \\
  n(\mathrm{H^+})_{\text{total}} &= 2\,n(\mathrm{H_2T})
     = 3.79\times10^{-2}\ \mathrm{mol}\\
  n(\mathrm{KMnO_4}) &\leq \frac{0.01 \times 2.878}{158.034}
     = 1.82\times10^{-4}\ \mathrm{mol}\\
  n(\mathrm{Zn})_{\text{mean}} &= \frac{1.266}{65.38}
     = 1.936\times10^{-2}\ \mathrm{mol}
     \quad (\text{range } 0.85\text{--}2.57\times10^{-2})\\
  n(\mathrm{HgBr_2}) &= \frac{1\times10^{-3}}{360.41}
     = 2.774\times10^{-6}\ \mathrm{mol}
\end{align}

Complete dissolution of the mean zinc charge would require
$2 \times 1.936\times10^{-2} = 3.87\times10^{-2}$~mol of H$^+$, against
$3.79\times10^{-2}$~mol available: the two are matched to within 2\%.
Across the observed zinc range, the acid provides between a 2.2-fold excess
(smallest sachet) and only 74\% of requirement (largest sachet).

Note also that the patent recommends $0.2$--$1$~g of zinc for a 100~mL
sample, and its worked examples use $0.29$--$0.70$~g of zinc against the
same $2.7$--$2.8$~g of acid. The sachets measured here average nearly twice
the top of that range. It is worth confirming that the kit is dosed for a
100~mL sample.

\subsection{Stoichiometric excess over arsenic}

At the WHO guideline value of 10~\ptl{} in a 100~mL sample:
\begin{equation}
  n(\mathrm{As}) = \frac{1.0\times10^{-6}}{74.92}
  = 1.34\times10^{-8}\ \mathrm{mol}
\end{equation}
Reduction of As(V) to arsine consumes 4~mol Zn per mol As, i.e.\
$5.3\times10^{-8}$~mol $= 3.5$~\ensuremath{\mu}g of zinc. The smallest
sachet contains 0.553~g. \textbf{Zinc is in approximately
$10^{4}$--$10^{5}$-fold excess and cannot be stoichiometrically limiting.}
The same conclusion holds for the acid.

\subsection{Indicator pad capacity}
\label{sec:padcap}

By~\eqref{eq:tri}, whole-pad capacity is
\begin{equation}
  n(\mathrm{As})_{\max} = \frac{n(\mathrm{HgBr_2})}{3}
  = 9.25\times10^{-7}\ \mathrm{mol} = 69.3\ \ensuremath{\mu}\mathrm{g}
\end{equation}
equivalent to 693~\ptl{} in a 100~mL sample \emph{if the arsenic were
distributed uniformly and captured completely}. The areal capacity of the
$32.5$~mm$^2$ pad is
\begin{equation}
  \sigma_{\text{cap}} = \frac{9.25\times10^{-7}}{32.5}
  = 28.5\ \mathrm{nmol\,mm^{-2}}
  \label{eq:sigmacap}
\end{equation}
This figure governs the onset of \emph{local} saturation in
Section~\ref{sec:deposition}, which occurs at far lower concentrations than
693~\ptl{}.

%======================================================================
\section{The real limiting factor: buffer capacity}
%======================================================================

\subsection{Derivation}

Since neither reagent is stoichiometrically limiting for arsenic, the
question becomes: how long does the system remain in the pH window where
arsine generation is fast?

\paragraph{Assumptions.}
(i)~Activity coefficients are taken as unity (ionic strength rises during
the run, so this is approximate).
(ii)~The only proton sink is zinc dissolution via~\eqref{eq:zn}; arsenic and
permanganate consumption are negligible.
(iii)~Sample alkalinity is neglected; a real groundwater with
$2$~meq\,L$^{-1}$ alkalinity consumes an additional $0.2$~mmol, roughly
5\% of the acid, and shifts the results accordingly.
(iv)~Zinc tartrate complexation is neglected; it will slightly increase the
apparent capacity.

Let $C_T$ be total tartrate, $V$ the sample volume, and define the
protonation fractions
\begin{equation}
  \alpha_1 = \frac{K_{a1}[\mathrm{H^+}]}{D}, \qquad
  \alpha_2 = \frac{K_{a1}K_{a2}}{D}, \qquad
  D = [\mathrm{H^+}]^2 + K_{a1}[\mathrm{H^+}] + K_{a1}K_{a2}
\end{equation}
Charge balance on the initial solution (pure H$_2$T, neglecting OH$^-$)
gives $[\mathrm{H^+}]_0 = C_T(\alpha_{1,0} + 2\alpha_{2,0})$, so the
proton condition referenced to the initial state reduces to
\begin{equation}
  \boxed{\;
  n_{\mathrm{H^+}}(\mathrm{pH}) = V\Big[
    C_T\big(\alpha_1 + 2\alpha_2\big) - [\mathrm{H^+}]
  \Big]\;}
  \label{eq:proton}
\end{equation}
which is the number of moles of protons that zinc has removed by the time
the solution reaches that pH. The corresponding dissolved zinc is
$n_{\mathrm{Zn}} = n_{\mathrm{H^+}}/2$.

\subsection{Result}

With $C_T = 0.190$~M, $V = 0.100$~L, $K_{a1} = 1.047\times10^{-3}$,
$K_{a2} = 4.571\times10^{-5}$:

\begin{table}[h]
\centering
\caption{Zinc dissolved before the solution reaches a given pH.}
\begin{tabular}{lrrr}
\toprule
pH & $n_{\mathrm{H^+}}$ (mmol) & $n_{\mathrm{Zn}}$ (mmol)
   & Zinc dissolved (g) \\
\midrule
1.87 (initial) & 0.00 & 0.00 & 0.000 \\
2.20 & 2.11 & 1.06 & 0.069 \\
2.50 & 4.53 & 2.26 & \textbf{0.148} \\
3.00 & 10.27 & 5.13 & 0.336 \\
\bottomrule
\end{tabular}
\end{table}

\paragraph{Interpretation.} Only about \textbf{0.15~g of zinc can dissolve
before the pH leaves the productive window} of pH~$\leq 2.5$. Even the
smallest sachet contains 3.7 times this amount. Zinc is therefore never
consumed as a reagent; it functions as a delivery vehicle for
\emph{surface area} and for the co-formulated rate catalyst. The acid's
buffer capacity is the true ceiling on how long the reaction runs hot, and
increasing the acid charge extends that window roughly proportionally.

\subsection{Hydrogen volume and the impossibility of a sealed vessel}

Complete consumption of the acid would liberate
$1.90\times10^{-2}$~mol of H$_2$, i.e.\ 465~mL at $25\,^\circ$C and 1~atm.
Reaching pH~2.5 alone liberates $2.26$~mmol $=55$~mL. Against a 73~mL
headspace, a gas-tight closure would reach
\begin{equation}
  P = \frac{nRT}{V} = \frac{1.90\times10^{-2}\times 8.314 \times 298}
  {7.3\times10^{-5}} \approx 6.4\ \mathrm{atm\ gauge}
\end{equation}
\textbf{The vessel must vent, and the venting path must be a designed,
reproducible feature rather than a leak.} An ill-fitting closure converts an
uncontrolled leak rate into an uncontrolled capture efficiency.

\subsection{Summary of limiting factors}

\begin{table}[h]
\centering
\caption{What limits what.}
\begin{tabular}{lll}
\toprule
Candidate & Limiting? & Basis \\
\midrule
Zinc vs.\ arsenic       & No, by $\sim\!10^5$ & Section 3.3 \\
Acid vs.\ arsenic       & No, by $\sim\!10^5$ & Section 3.3 \\
Acid \emph{buffer capacity} & \textbf{Yes} & Table 2 \\
HgBr$_2$, whole pad     & Only above $\sim\!690$~\ptl{} & Section 3.4 \\
HgBr$_2$, \emph{locally}& \textbf{Yes, from tens of \ptl{}}
                                              & Section 6.4 \\
Reaction time           & \textbf{Yes} & Section 5 \\
Arsine capture efficiency & \textbf{Yes, and flow-dependent}
                                              & Section 6 \\
\bottomrule
\end{tabular}
\end{table}

%======================================================================
\section{Kinetics, timing, and temperature}
%======================================================================

\subsection{Reading before the plateau}

The patent states that substantially complete arsenic reduction may occur
within about 10 minutes, and all of its worked examples read at 10 minutes.
The kit protocol reads at 7 minutes. Reading on the rising part of the
conversion curve is the single most consequential procedural choice, because
it converts every rate variable --- zinc mass, catalyst dose, particle size,
temperature, pH trajectory --- directly into a concentration error.

The sensitivity of the response to zinc and catalyst loading is documented
in the patent's own examples. Against a fixed 100~ppb standard:

\begin{table}[h]
\centering
\caption{Patent examples: apparent concentration for a 100~ppb standard.}
\begin{tabular}{llll}
\toprule
Example & Zinc reagent & Zinc (g) & Reported result \\
\midrule
11 & T2 & 0.29 & $>100$, $<150$~ppb \\
10 & T1 & 0.70 & $>200$, $<250$~ppb \\
14 & T1 & 0.70 (higher KMnO$_4$) & $\approx 400$~ppb \\
\bottomrule
\end{tabular}
\end{table}

\noindent
A 2.4-fold change in zinc and catalyst moved a fixed sample across a
fourfold range of apparent concentration. The sachets measured here vary
threefold.

\subsection{Time--temperature compensation}

\paragraph{Assumptions.}
(i)~A single activated process is rate-controlling.
(ii)~The reaction proceeds to the \emph{same endpoint} at all temperatures,
so that only the pace differs. This must be verified
(Section~\ref{sec:isoconv-check}); it can fail if a hotter run exhausts the
buffer window before conversion is complete.
(iii)~The Arrhenius form applies over the narrow range of interest.

Isoconversion requires $k(T)\,t = \text{constant}$. With
$k = A\exp(-E_a/RT)$,
\begin{equation}
  \boxed{\;
  t(T) = t_{25}\,\exp\!\left[
    \frac{E_a}{R}\left(\frac{1}{T} - \frac{1}{298.15}\right)
  \right]\;}
  \label{eq:arrhenius}
\end{equation}
with $T$ in kelvin, $R = 8.314$~J\,mol$^{-1}$\,K$^{-1}$, and $t_{25}$ the
chosen reference reaction time. This holds regardless of reaction order.

It is convenient to recast $E_a$ as a $Q_{10}$, the rate factor per
$10\,^\circ$C:
\begin{equation}
  Q_{10} = \exp\!\left[\frac{E_a}{R}\cdot
  \frac{10}{T(T+10)}\right], \qquad
  t(\theta) = t_{25}\,Q_{10}^{\,(25-\theta)/10}
  \label{eq:q10}
\end{equation}
At $T = 298$~K, $Q_{10} = 2$ corresponds to
$E_a \approx 53$~kJ\,mol$^{-1}$, which lies in the middle of the published
range for surface-controlled acid dissolution of zinc
($\approx 40$--$60$~kJ\,mol$^{-1}$; diffusion control would give
$\approx 20$).

\begin{table}[h]
\centering
\caption{Compensated read times for $t_{25} = 12$~min, $Q_{10} = 2$.
Use only after measuring your own $Q_{10}$.}
\begin{tabular}{lrrrrrrrrrrr}
\toprule
Water $^\circ$C & 15 & 18 & 20 & 22 & 24 & 25 & 26 & 28 & 30 & 32 & 35\\
\midrule
Read at (min) & 24.0 & 19.5 & 17.0 & 14.8 & 12.9 & 12.0 & 11.2
              & 9.7 & 8.5 & 7.4 & 6.0 \\
\bottomrule
\end{tabular}
\end{table}

If temperature drifts during a run, accumulate equivalent time and stop when
it reaches $t_{25}$:
\begin{equation}
  t_{\mathrm{eq}} = \int_0^{t}\exp\!\left[
    \frac{E_a}{R}\left(\frac{1}{298.15} - \frac{1}{T(\tau)}\right)
  \right]\mathrm{d}\tau
  \label{eq:integral}
\end{equation}

\subsection{Measuring $E_a$ empirically}

Run a fixed standard ($\approx 50$~\ptl{}) at two temperatures as far apart
as ambient allows, with a full time course at each (5, 8, 12, 18, 25, 35~min,
in triplicate). Find the times $t_1, t_2$ giving \emph{equal} colour:
\begin{equation}
  E_a = \frac{R\,\ln(t_1/t_2)}{\dfrac{1}{T_1} - \dfrac{1}{T_2}}
  \qquad\text{or}\qquad
  Q_{10} = \left(\frac{t_1}{t_2}\right)^{\frac{10}{\theta_2-\theta_1}}
\end{equation}

The rigorous form of the same experiment is time--temperature superposition:
plot response against $\ln t$ for each temperature and shift the curves
horizontally until they superimpose. The required shift is $\ln a_T$. If the
curves collapse onto a single master curve, the single-process assumption is
validated; if they do not, no time correction will work.

\subsection{Verifying the isoconversion assumption}
\label{sec:isoconv-check}

Compare the \emph{plateau heights} at the two temperatures.
\begin{itemize}[nosep]
  \item Plateaus match $\Rightarrow$ adjust time
        via~\eqref{eq:arrhenius}.
  \item Plateaus differ $\Rightarrow$ the endpoint itself is
        temperature-dependent, which is what one expects if the buffer
        window is exhausted before conversion completes. A multiplicative
        correction on the reading, $R_{\mathrm{corr}}
        = R_{\mathrm{meas}}\,f(\theta)$, must then be fitted empirically.
        Raising Reagent~1 to $5$--$6$~g per 100~mL widens the buffer window
        and should be tried first.
\end{itemize}

\subsection{Choosing a time where compensation barely matters}

The better strategy is to operate where the response is flat. A temperature
uncertainty of $\pm\Delta\theta$ at a given $Q_{10}$ produces an effective
time uncertainty
\begin{equation}
  \Delta\ln t = \frac{\Delta\theta}{10}\ln Q_{10}
\end{equation}
The resulting relative error in the reading is
\begin{equation}
  \frac{\Delta R}{R} = \left(\frac{\mathrm{d}\ln R}{\mathrm{d}\ln t}\right)
  \Delta \ln t
\end{equation}
For $\Delta\theta = 5\,^\circ$C and $Q_{10} = 2$,
$\Delta\ln t = 0.347$. Requiring $\Delta R/R \leq 0.02$ gives the
\emph{flatness criterion}
\begin{equation}
  \boxed{\;\frac{\mathrm{d}\ln R}{\mathrm{d}\ln t} \leq 0.058\;}
\end{equation}
which converts to a one-line bench test. Read the same standard at 10~min
and at 20~min ($\Delta\ln t = \ln 2 = 0.693$):
\begin{equation}
  \frac{R_{20} - R_{10}}{R_{10}} \leq 0.058 \times 0.693 = 0.040
\end{equation}
\textbf{If the 10-minute and 20-minute readings differ by 4\% or less, a
$\pm 5\,^\circ$C temperature excursion costs under 2\% and the correction
can be omitted.} If they differ by 15\%, the correction is doing real work
and $t_{25}$ should be pushed toward 15~min.

\subsection{Thermal budget}

The sample is not isothermal. For~\eqref{eq:zn},
$\Delta H \approx -154$~kJ\,mol$^{-1}$. Dissolution of 0.3~g of zinc
(4.59~mmol) into 100~g of water ($c_p = 4.18$~J\,g$^{-1}$K$^{-1}$) releases
\begin{equation}
  \Delta T = \frac{4.59\times10^{-3}\times 1.54\times10^{5}}
  {100 \times 4.18} = +1.7\ \mathrm{K}
\end{equation}
Dissolution of tartaric acid is endothermic
($\Delta H_{\mathrm{sol}} \approx +20$ to $+25$~kJ\,mol$^{-1}$), giving
roughly $-1$~K during Step~1. \textbf{Record the temperature at the moment
Reagent~2 is added and again at the read, and use the mean
in~\eqref{eq:arrhenius}.} Ignoring a $1.5$~K rise costs about 10\% of the
effective time.

This also has a bearing on any strategy of pre-warming the pad holder to
suppress condensation: the water warms during the run while a small plastic
head loses heat to the room, so the temperature margin closes exactly when
it is most needed. Start the head $5$--$8\,^\circ$C above the sample and
insulate the tube above the closure.

%======================================================================
\section{Gas handling and pad exposure geometry}
%======================================================================

\subsection{What the geometry must achieve}

Colour development is more efficient if the gas is forced into intimate
contact with the reagent paper and if the exposed area is small enough to
concentrate the available arsine. For a pad bonded to an impermeable
carrier, literal through-flow is impossible, but the two functional
requirements can still be met:
\begin{enumerate}[nosep]
  \item \textbf{No bypass} --- every molecule of carrier gas must traverse
        the pad face. The back face must be sealed.
  \item \textbf{Short diffusion distance} --- the gap over the pad face must
        be small enough that arsine reaches the surface within the residence
        time.
\end{enumerate}
Failure of requirement~1 was observed directly: with the pad merely
suspended in the vial airspace, colour developed markedly more strongly
around the edges, which is the signature of the pad mapping an
uncontrolled flow field rather than integrating a captured dose.

\subsection{Radial flow between the pad and a flat face}

Consider gas admitted through a central hole of radius $r_1$ into a gap of
height $h$ between the pad face and a flat plate, escaping at radius $r_2$.

\paragraph{Assumptions.} Creeping (Stokes) flow; $h \ll r$; isothermal,
incompressible, Newtonian; the pad face is smooth on the scale of $h$;
entrance effects neglected.

Local volumetric flow through the annulus at radius $r$:
\begin{equation}
  Q = -(2\pi r)\frac{h^3}{12\mu}\frac{\mathrm{d}P}{\mathrm{d}r}
\end{equation}
Separating and integrating from $r_1$ to $r_2$:
\begin{equation}
  \boxed{\;
  Q = \frac{\pi h^3 \Delta P}{6\mu \ln(r_2/r_1)}
  \quad\Longleftrightarrow\quad
  h = \left[\frac{6\mu Q \ln(r_2/r_1)}{\pi \Delta P}\right]^{1/3}\;}
  \label{eq:radialflow}
\end{equation}

\subsection{The force-balanced (flapper valve) case}

If the pad is held against the flat face by a fixed light force $F$, the gap
self-adjusts until the integrated pressure lifts the pad. The pressure
profile from~\eqref{eq:radialflow} is
\begin{equation}
  P(r) - P_{\mathrm{atm}} = \Delta P\,
  \frac{\ln(r_2/r)}{\ln(r_2/r_1)}
\end{equation}
Integrating over the annulus and adding the hole area, with
$L \equiv \ln(r_2/r_1)$:
\begin{equation}
  F = \Delta P\left[\frac{2\pi}{L}
  \left(\frac{r_2^2 - r_1^2}{4} - \frac{r_1^2 L}{2}\right)
  + \pi r_1^2\right]
  = \Delta P\,\underbrace{\frac{\pi(r_2^2-r_1^2)}{2L}}_{A_{\mathrm{eff}}}
  \label{eq:aeff}
\end{equation}
The $r_1^2$ terms cancel exactly, leaving a compact effective area.

For a 3/32~in hole ($r_1 = 1.19$~mm) and a 5~mm-wide pad
($r_2 = 2.5$~mm to the nearest edge), $L = 0.742$ and
$A_{\mathrm{eff}} = 10.2$~mm$^2$. A closing force of $0.5$~gf gives
$\Delta P = 479$~Pa. With $\mu(\mathrm{H_2}, 25\,^\circ\mathrm{C})
= 8.9\times10^{-6}$~Pa\,s and a representative peak gas rate
$Q = 0.9$~mL\,s$^{-1}$, equation~\eqref{eq:radialflow} gives
\begin{equation}
  h \approx 29\ \um
\end{equation}

\paragraph{Key scaling.} Because $\Delta P$ is fixed by the spring force and
is \emph{independent of flow},
\begin{equation}
  h \propto Q^{1/3}
\end{equation}
A doubling of gas rate opens the gap by only 26\%. \textbf{The intuition
that a hotter, faster reaction would exponentially widen the gap is
incorrect}; the dependence is a cube root. Temperature still matters
enormously, but through the generation chemistry of Section~5, not through
the valve.

\subsection{Capture efficiency}

\paragraph{Assumptions.} Plug flow in the radial direction; wall mass
transfer coefficient $k_m \approx D/h$; arsine concentration uniform across
the gap; the pad is a perfect sink (reaction fast relative to diffusion);
one absorbing wall.

Differential mass balance, with $\mathrm{d}A$ the incremental pad area:
\begin{equation}
  Q\,\mathrm{d}C = -k_m C\,\mathrm{d}A
  \quad\Longrightarrow\quad
  \frac{C(A)}{C_0} = \exp\!\left(-\frac{A}{A^{*}}\right),
  \qquad
  \boxed{\;A^{*} = \frac{hQ}{D}\;}
  \label{eq:astar}
\end{equation}
$A^{*}$ is the \emph{e-folding area}: the pad area over which 63\% of the
arsine is removed. The number of transfer units for the whole accessible
pad area $A_p$ is
\begin{equation}
  \mathrm{NTU} = \frac{A_p}{A^{*}} = \frac{A_p D}{hQ},
  \qquad
  \eta_{\text{capture}} = 1 - e^{-\mathrm{NTU}}
  \label{eq:ntu}
\end{equation}

The assumption $k_m = D/h$ is deliberately conservative: for fully developed
flow in a slot with one absorbing wall, $\mathrm{Sh} \approx 4.86$ gives
$k_m \approx 2.4\,D/h$, which would raise all NTU values by that factor.

Taking $D(\mathrm{AsH_3}$ in $\mathrm{H_2}) \approx
4\times10^{-5}$~m$^2$\,s$^{-1}$ (estimated; the value is uncertain within
roughly $3$--$7\times10^{-5}$) and $A_p \approx 25$~mm$^2$:

\begin{table}[h]
\centering
\caption{Force-balanced valve: $\mathrm{NTU} \propto Q^{-4/3}$
because $h \propto Q^{1/3}$.}
\begin{tabular}{lrrrr}
\toprule
$Q$ (mL\,s$^{-1}$) & $h$ (\um) & $A^{*}$ (mm$^2$) & NTU
& $\eta_{\text{capture}}$ \\
\midrule
0.3 & 19.9 & 0.149 & 168 & $>0.9999$ \\
0.9 & 28.7 & 0.646 & 38.7 & $>0.9999$ \\
2.7 & 41.4 & 2.79 & 8.9 & 0.9999 \\
8.1 & 59.7 & 12.1 & 2.1 & 0.876 \\
\bottomrule
\end{tabular}
\end{table}

\noindent
Capture is effectively quantitative across the full range of gas rates the
threefold zinc variation can produce. \textbf{This is the decisive advantage
of a sealed, back-supported geometry: it converts collection from a
flow-dependent efficiency into a near-quantitative trap, and thereby
decouples the reading from the zinc mass.}

\subsection{Deposition profile and local saturation}
\label{sec:deposition}

The same analysis fixes \emph{where} the arsenic lands. Areal deposition
density for a total captured amount $n$:
\begin{equation}
  \sigma(A) = \frac{n}{A^{*}}\exp\!\left(-\frac{A}{A^{*}}\right)
  \label{eq:sigma}
\end{equation}
peaking at the hole edge. Local saturation begins when
$\sigma(0) = \sigma_{\text{cap}}$ from~\eqref{eq:sigmacap}:
\begin{equation}
  \boxed{\;n_0 = A^{*}\sigma_{\text{cap}}\;}
  \label{eq:onset}
\end{equation}
and above that, the saturated area grows only logarithmically:
\begin{equation}
  A_{\text{sat}}(n) = A^{*}\ln\!\left(\frac{n}{n_0}\right)
  \label{eq:asat}
\end{equation}

\paragraph{Consequence for the flapper valve.} With
$A^{*} = 0.646$~mm$^2$, equation~\eqref{eq:onset} gives
$n_0 = 18.4$~nmol $= 1.38$~\ensuremath{\mu}g, i.e.\ saturation onset at
about \textbf{14~\ptl{}}. Above that the stain does not darken further at
its centre; it widens --- and by~\eqref{eq:asat} the widening is
logarithmic, so between 14 and 256~\ptl{} the saturated radius grows only
from 1.19 to about 1.46~mm. The integrated response is strongly sublinear
over almost the entire useful range.

\paragraph{The fundamental trade-off.} Combining~\eqref{eq:ntu}
and~\eqref{eq:onset}:
\begin{equation}
  \mathrm{NTU} = \frac{A_p}{A^{*}}
  \qquad\text{and}\qquad
  n_0 \propto A^{*}
\end{equation}
\textbf{Quantitative capture and uniform deposition are in direct
opposition.} Large NTU (good capture) requires $A^{*} \ll A_p$, which
concentrates the deposit near the inlet. Setting $A^{*} \approx A_p$ for
uniform deposition drops capture to 63\%. There is no geometry that gives
both, and the design task is to choose the compromise deliberately.

\subsection{Recommended geometry: a fixed shim-defined gap}
\label{sec:fixedgap}

Since $A^{*} = hQ/D$, the gap height is the single lever controlling both
capture and linear range. Requiring $\mathrm{NTU} \geq 3$ (i.e.\
$\eta \geq 95$\%) at the maximum expected gas rate gives the design rule
\begin{equation}
  \boxed{\;h \leq \frac{A_p D}{3\,Q_{\max}}\;}
  \label{eq:designrule}
\end{equation}
With $A_p = 25$~mm$^2$, $D = 4\times10^{-5}$~m$^2$\,s$^{-1}$ and
$Q_{\max} = 1.5$~mL\,s$^{-1}$, this gives $h \leq 222$~\um{}. A gap of
$150$~\um{} leaves margin:

\begin{table}[h]
\centering
\caption{Fixed gap $h = 150$~\um{}: $\mathrm{NTU} \propto Q^{-1}$.}
\begin{tabular}{lrrrr}
\toprule
$Q$ (mL\,s$^{-1}$) & $A^{*}$ (mm$^2$) & NTU
& $\eta_{\text{capture}}$ & Saturation onset $n_0$ \\
\midrule
0.3 & 1.13 & 22.2 & $>0.9999$ & 32~nmol (24~\ptl{})\\
0.9 & 3.38 & 7.4  & 0.9994 & 96~nmol (72~\ptl{})\\
1.5 & 5.63 & 4.4  & 0.988  & 160~nmol (120~\ptl{})\\
\bottomrule
\end{tabular}
\end{table}

\noindent
Compared with the $29$~\um{} force-balanced gap, a fixed $150$~\um{} gap
gives:
\begin{itemize}[nosep]
  \item capture still $\geq 98.8$\% across the whole plausible flow range;
  \item a fivefold extension of the linear range
        (14~\ptl{} $\rightarrow$ 72~\ptl{} onset at nominal flow);
  \item \textbf{no force to calibrate} --- the gap is set by a shim, not by
        a spring balance at the milligram-force scale;
  \item \textbf{no valve chatter} at the tail of the run, when gas
        generation falls below the cracking pressure of a force-balanced
        design;
  \item \textbf{no capillary stiction} --- the capillary pressure of a
        condensed water bridge is $2\gamma/h$, which falls from
        $\approx 5000$~Pa at 29~\um{} to $\approx 960$~Pa at 150~\um{}, and
        in any case there is no closing force for a bridge to act against;
  \item \textbf{no temperature sensitivity of the gap at all}, since $h$ is
        a machined dimension.
\end{itemize}

\paragraph{Construction.} A rectangular channel a little wider than the
strip ($\approx 5.3$~mm), with the strip's plastic back seated flat against
the channel floor and a shim or milled step establishing $150$~\um{} over
the pad face. Gas enters through a transverse slot ($\approx 5 \times
0.3$~mm) rather than discrete round holes, so that it arrives as a sheet
spanning the pad width. Extend the channel $10$--$15$~mm beyond the pad
before the exit, and place any flow restriction \emph{downstream} of the
pad. Do not use a second floating pressure regulator (such as an inverted
vial) in series; two floating elements can hunt against one another. Vent
through a fixed open orifice and make the closure itself genuinely
gas-tight (silicone or butyl rather than natural cork).

\paragraph{Note on rough, porous pad surfaces.} Paper pore structure is
typically $1$--$20$~\um{}. At a $150$~\um{} working gap the paper's
roughness is a small perturbation on the channel height, and the paper wicks
condensate into its bulk rather than allowing a coherent liquid bridge to
span the gap. This is a genuine advantage of paper over a smooth absorber,
though wicking capacity is finite: pale, diffuse, soft-edged stains rather
than a crisp pattern are the signature of a flooded pad.

\subsection{The bypass failure mode}

Whatever the gap, the claimed advantage --- that all gas passes close to the
pad --- is destroyed by any preferential leak path: a burr at the inlet, a
scratch across the flat face, a wave in the plate, or a domed fillet at a
resin edge. Such a defect is invisible on inspection but not in the data.

\textbf{The stain shape is a free diagnostic.} A clean concentric or
transversely symmetric pattern indicates a good seal. A crescent, comet
tail, or off-centre blob indicates tilt or a leak. Photograph every pad and
treat asymmetry as grounds for rejecting the run.

\subsection{Instrumenting the pressure}

Rather than attempting to calibrate a milligram-scale closing force, tee a
length of clear tubing partly filled with water off the reaction vessel as a
water manometer, with a millimetre scale behind it. This gives, live and for
the cost of a metre of tubing: the operating $\Delta P$, run-to-run
repeatability of the gas path, visible evidence of any chatter, and a
qualitative trace of the gas generation curve.

\subsection{Two decisive experiments}

\paragraph{Tandem pad breakthrough.} Place a second pad downstream of the
first and run a 100~\ptl{} standard. If the capture analysis is correct,
the downstream pad should remain white. The ratio of downstream to upstream
colour is a direct measurement of $1 - \eta$, and hence a check on the
assumed $D$ and $k_m$. Repeat at 0.7~g and 1.6~g of zinc: if the downstream
pad stays clean at both, the geometry has decoupled the reading from the gas
rate, and remaining error must lie in the chemistry.

\paragraph{Gas volumetry.} Route the outlet through tubing to an inverted
graduated cylinder over water and record volume against time. This measures
the hydrogen evolution curve directly, giving $Q(t)$, $Q_{\max}$ for use
in~\eqref{eq:designrule}, and a direct test of whether duplicate runs at
identical zinc mass produce superimposable curves.

%======================================================================
\section{Reflectance photometry}
\label{sec:optics}
%======================================================================

\subsection{Kubelka--Munk}

Diffuse reflectance from a scattering paper substrate is not linear in
absorber concentration. For an optically thick, diffusely scattering layer,
Kubelka--Munk theory gives
\begin{equation}
  \boxed{\;
  f(R_\infty) = \frac{(1 - R_\infty)^2}{2R_\infty}
  = \frac{K}{S}\;}
  \label{eq:km}
\end{equation}
where $K$ is the absorption coefficient and $S$ the scattering coefficient.

\paragraph{Assumptions.}
(i)~The layer is optically thick (no light reaches the backing) --- verify
by measuring a pad over white and over black backings and confirming the
readings agree.
(ii)~Isotropic diffuse illumination and collection; a directional LED and
camera approximate this only roughly, so treat $f(R)$ as an empirical
linearising transform rather than an absolute quantity.
(iii)~$S$ is independent of wavelength and of the absorber --- reasonable
for a lightly loaded pad, but violated by wetting.
(iv)~$R_\infty$ is measured relative to a stable diffuse white standard.

Because $K$ is proportional to chromophore concentration while $S$ is a
property of the paper, subtract the unexposed pad measured in the same band:
\begin{equation}
  \Delta(K/S)_\lambda = f(R_\lambda)_{\text{exposed}}
  - f(R_\lambda)_{\text{blank}}
  \label{eq:dks}
\end{equation}
Only $\Delta(K/S)$ is approximately proportional to arsenic loading.
Performing linear algebra on raw reflectance, or on gamma-encoded camera
values, manufactures curvature that is not chemistry.

\subsection{What the spectra should look like}

The three substitution states are not a freely variable mixture: they are
sequential rungs whose relative abundance is set by the \emph{local}
As:Hg ratio. Resolving them therefore yields the local saturation state, not
an independent measure of total arsenic. The correct framing is two
observables from one spectrum:
\begin{itemize}[nosep]
  \item \textbf{magnitude} of $\Delta(K/S)$ --- how much arsenic arrived;
  \item \textbf{shape} of $\Delta(K/S)$ --- how concentrated it was where it
        landed.
\end{itemize}

Do not expect three resolvable peaks. These are charge-transfer absorptions
in a progressively more arsenide-like solid, so advancing substitution
produces a \emph{red-shifting absorption edge} rather than new bands.
``Yellow'' is an edge near $450$--$470$~nm; ``brown'' is that edge dragged
toward $550$--$600$~nm with an extended tail.

\subsection{LED band selection}

Bands should straddle the absorption edge, not be spread evenly across the
visible.

\begin{table}[h]
\centering
\caption{Suggested narrow-band LED set (4-band minimum; 8-band preferred).}
\begin{tabular}{llll}
\toprule
$\lambda$ (nm) & Priority & Dominant Bayer channel & Purpose \\
\midrule
415--430 & Essential & B & Core of the yellow absorption \\
450--470 & Essential & B & Edge position, mono/di-substituted \\
495--505 & Preferred & B/G & Edge shoulder \\
525--535 & Essential & G & Brown/trisubstituted region \\
565--590 & Preferred & G/R & Far edge tail; neutral absorbers \\
620--630 & Essential & R & Reference: no arsenic species absorbs \\
655--670 & Preferred & R & Second reference; scattering baseline \\
730--780 & Optional  & R & Pure scattering; wetting and gloss \\
\bottomrule
\end{tabular}
\end{table}

\paragraph{Practical notes.}
Use direct-emission narrow-band LEDs, never phosphor-converted white types.
Avoid wavelengths below about $410$~nm: papers frequently contain optical
brighteners that fluoresce blue under violet and UV excitation, which would
appear as a spurious, dose-independent signal --- test a blank pad at
$405$~nm before committing to any short-wavelength band. Characterise each
LED's peak and FWHM rather than trusting the datasheet, since a $20$~nm
centroid error near the absorption edge is a large photometric error.
Correct for Bayer-channel crosstalk by measuring each LED's response in all
three channels on a neutral target.

\subsection{Establishing whether the shape carries information}

Run a dose series (0, 8, 16, 32, 64, 128, 256~\ptl{}, triplicate) with
everything else fixed and collect all bands.

\paragraph{Cheap test.} Normalise each $\Delta(K/S)$ spectrum to unit sum
and overlay. If they superimpose, only the amount is varying and
multispectral analysis adds nothing beyond a better single number. If they
fan out systematically with dose, composition is changing. Look also for an
\textbf{isosbestic point}, a wavelength at which all normalised curves
cross: its presence is direct evidence of two species interconverting in
fixed stoichiometry.

\paragraph{Rigorous test.} Assemble the $\Delta(K/S)$ values into a matrix
$\mathbf{X}$ (samples $\times$ bands) and compute the singular value
decomposition $\mathbf{X} = \mathbf{U}\boldsymbol{\Sigma}\mathbf{V}^{\!\top}$.
The number of singular values rising above the noise floor is the number of
independently varying chromophores. If $\sigma_2/\sigma_1$ sits at
instrument noise, the system is rank~1. If rank $\geq 2$, apply
multivariate curve resolution with alternating least squares (MCR-ALS) under
non-negativity constraints on both spectra and concentrations, which
recovers the pure component spectra without isolated standards.

\paragraph{Interpreting rank 1.} Rank~1 does not prove a single species; it
proves the species \emph{ratio} is constant across samples. That is
excellent news for calibration --- a single well-behaved parameter --- while
indicating that multispectral shape analysis adds nothing.

\subsection{The spatial requirement}

An area-averaged spectrum of a saturating stain cannot distinguish
``uniformly half-substituted'' from ``core saturated, periphery untouched''
--- these give the same integrated spectrum at very different loadings.
Since the whole purpose of the shape metric is to detect saturation, and
saturation is spatially localised by~\eqref{eq:sigma}, \textbf{the spectrum
must be evaluated per-pixel and then integrated}, not averaged first.

This is exactly what a camera-based instrument can do and a single-aperture
colorimeter cannot. The recommended readout is:
\begin{align}
  \text{low range:}\quad & \Sigma = \sum_{\text{pixels}}
    \Delta(K/S)_{\lambda_{\text{blue}}} \cdot \delta A
    \quad\text{(linear below } n_0) \\
  \text{high range:}\quad & A_{\text{sat}} =
    \text{area of pixels classified as saturated}
    \quad\text{(logarithmic, per~\eqref{eq:asat})}
\end{align}
spliced at the onset $n_0$ determined empirically.

\subsection{Correction for pad wetting}

This is arguably the most valuable single application of the multi-band
instrument. Water darkens paper \emph{neutrally}, by changing the scattering
coefficient $S$ across all wavelengths. Arsenic darkens \emph{selectively}
in the blue, by changing $K$. A reference band at $620$--$670$~nm therefore
measures wetting almost independently of arsenic loading and permits its
removal:
\begin{equation}
  \Delta(K/S)^{\text{corr}}_\lambda
  = \Delta(K/S)_\lambda
  - \beta_\lambda\,\Delta(K/S)_{\text{ref}}
\end{equation}
where $\beta_\lambda$ is calibrated from deliberately wetted blank pads.
A single-band instrument cannot separate these effects at all; it reads a
damp blank as arsenic. Given the condensation risk inherent in a
near-saturated gas stream, this correction should be regarded as mandatory.

\subsection{Cautions}

\begin{itemize}[nosep]
  \item \textbf{Neutral absorbers masquerade as brown.} Hg$_3$As$_2$ and
        metallic deposits are spectrally flat, and a flat added absorbance
        moves any blue/red ratio metric in the same direction as genuine
        trisubstitution. The flat-versus-edge distinction is precisely what
        the red reference band resolves.
  \item \textbf{Sulfide is not a clean grey flag.} Although HgS is black,
        the standard mercuric bromide paper documentation warns that
        hydrogen sulfide can itself produce a yellow colour. Test directly
        with a 0~\ptl{} sample spiked with 2~mg\,L$^{-1}$ sulfide and
        determine whether its spectral shape separates from a genuine
        arsenic stain.
  \item \textbf{Fading.} The yellow product fades, especially in sunlight,
        and there is no reason the three species should fade at equal rates.
        Measure one pad repeatedly over 20~min and check whether any ratio
        metric drifts independently of magnitude. Fix the interval between
        removal and measurement.
  \item \textbf{Optical thickness.} Confirm the pad is optically thick
        before relying on~\eqref{eq:km}.
\end{itemize}

%======================================================================
\section{Error model and quality assurance}
%======================================================================

Precision should be characterised over a range of concentrations rather than
at a single point. The standard approach is the Thompson--Howarth model, in
which the standard deviation of replicates varies with concentration as
\begin{equation}
  s_c = s_0 + k\,c
  \label{eq:th}
\end{equation}
where $s_0$ is the standard deviation at zero concentration (so that the
detection limit is $3s_0$) and $k$ is the constant relative error. Both
parameters can be estimated from duplicate analyses collected routinely
during normal work.

Published benchmarks for context: various laboratory procedures achieve
$s_0 = 0.30$~\ptl{}, $k = 0.088$; a hand-held instrumental Gutzeit device
gave $s_0 = 1.4$~\ptl{}, $k = 0.14$; visual field kits fall in the range
$s_0 = 5$--$7$~\ptl{}, $k = 0.2$--$0.3$. Instrumental colour detection can
in principle reach detection limits of $0.5$~\ptl{} for transmission and
$1$--$2$~\ptl{} for reflectance geometries.

Bias is a separate matter from precision and must be checked against
standards rather than against the supplied colour chart: an independent
evaluation of eight commercial kits found that one widely used kit was
repeatable between replicates but overestimated arsenic by roughly a factor
of two.

\paragraph{Routine QA.} Each session: one reagent blank; standards at 0, 10,
25, 50 and 100~\ptl{}; duplicates on 10\% of samples. Track $s_0$ and $k$
over time as the primary process-control metric. Verify As(III) and As(V)
standards separately, since their hydride-formation kinetics differ.

%======================================================================
\section{Revised protocol suggestions}
%======================================================================

\subsection{Changes requiring no recalibration}

\begin{enumerate}[nosep]
  \item \textbf{Restore the 2-minute oxidation rest} (patent range: 2 to
        5~min). Better, wait until the permanganate colour has visibly
        discharged, which makes the step self-indicating and removes the
        speciation variable.
  \item \textbf{Confirm complete dissolution} of Reagent~1 before adding
        Reagent~2; undissolved acid is unavailable acid.
  \item \textbf{Never seal the vessel}; vent through a fixed orifice
        downstream of the pad, and use a gas-tight elastomeric closure
        rather than porous cork.
  \item \textbf{Fix the stirring regime}: same follower, same rpm,
        continuous through the run rather than stirring then settling.
  \item \textbf{Fix the read interval}: pad out, measured within a set
        window, shielded from ambient light.
  \item \textbf{Record temperature twice} (at Reagent~2 addition and at
        read) and use the mean.
  \item \textbf{Photograph every pad} and reject asymmetric stains.
\end{enumerate}

\subsection{Reagent dispensing}

Pool approximately 20 sachets of each reagent, tumble-mix thoroughly, store
over desiccant, and dispense by mass: Reagent~1 at $2.90 \pm 0.02$~g,
Reagent~2 at $1.30 \pm 0.02$~g. This converts a 20\% CV input into
$\approx 1.5$\%.

\emph{Cautions.} Zinc dust is a flammable solid and degrades on humid
storage. If the catalyst salt differs in particle size from the zinc it may
segregate on pooling; verify with standards before and after pooling. Store
the permanganate-bearing Reagent~1 away from the zinc: an oxidiser and a
fine metal powder should never share a container.

\subsection{Finding a plateau operating point}

Run these in triplicate against a 50~\ptl{} As(V) standard:
\begin{enumerate}[nosep]
  \item \textbf{Time course}: 5, 7, 10, 15, 20, 30~min. Read where the
        curve flattens; apply the 4\% criterion of Section~5.6.
  \item \textbf{Zinc series}: 0.4, 0.7, 1.0, 1.3, 1.7~g at fixed acid and
        the new read time. Choose a point on the plateau.
  \item \textbf{Acid series}: 2.9, 4.0, 5.0, 6.0~g at fixed zinc. More acid
        widens the buffer window (Table~2) and should move the time-course
        plateau earlier and flatten the zinc dependence.
\end{enumerate}
The expected landing point is roughly $5$--$6$~g of Reagent~1,
$0.7$--$1.0$~g of Reagent~2, and a $15$--$20$~min read, trading throughput
for reproducibility. Any of these changes voids the manufacturer's colour
calibration, which must be rebuilt against standards regardless.

\subsection{Geometry}

Adopt the fixed shim-defined gap of Section~\ref{sec:fixedgap}:
back-sealed pad, $\approx 150$~\um{} channel over the pad face, transverse
slot inlet, restriction downstream, single venting element, water manometer
on the vessel. Validate with the tandem-pad breakthrough test and gas
volumetry.

\subsection{Optical}

Build the $\Delta(K/S)$ pipeline of Section~\ref{sec:optics} with per-pixel
evaluation, a red reference band for wetting correction, and the spliced
low-range/high-range readout. Establish the spectral rank before investing
in curve resolution.

%======================================================================
\section{Safety}
%======================================================================

\subsection{Arsine}

Occupational exposure limits: ACGIH TLV-TWA $0.005$~ppm ($\approx
16$~\ensuremath{\mu}g\,m$^{-3}$) as an 8-hour time-weighted average; OSHA
PEL $0.05$~ppm; NIOSH REL $0.002$~mg\,m$^{-3}$ as a 15-minute ceiling;
NIOSH IDLH $3$~ppm.

A 256~\ptl{} test releases at most $\approx 27$~\ensuremath{\mu}g of
arsine. At 30~CFM ($0.85$~m$^3$\,min$^{-1}$) over a 7-minute run, complete
release into $\approx 6$~m$^3$ of air gives $\approx 4.4$~\ensuremath{\mu}g
\,m$^{-3}$, below the TLV but without a large margin. The margin shrinks
proportionally with higher standards and with tests run back to back.

A computer fan in a 6~in duct realistically moves $20$--$40$~CFM. Through a
$12 \times 8$~in opening that is roughly $0.2$~m\,s$^{-1}$ ($\approx
45$~fpm) face velocity, against the $0.4$--$0.5$~m\,s$^{-1}$ expected of a
laboratory fume hood. A smoke test confirms direction, not velocity.
Recommended: measure face velocity with a hot-wire anemometer and target
$\geq 0.4$~m\,s$^{-1}$; keep the working aperture small; cap standards at
500~\ptl{} and dilute above that, which also addresses pad saturation.

\subsection{Hydrogen}

Maximum $465$~mL per run. Diluted into $\approx 6$~m$^3$ of ventilation air
this is $0.008$\%, two and a half orders of magnitude below the 4\% lower
explosive limit; ignition risk at the fan is negligible at that dilution.
Do not operate in a closed cabinet, and keep the reaction vessel away from
the fan.

\subsection{Other}

Reagent~1 contains an oxidiser; Reagent~2 is a flammable metal powder;
store separately. Used pads are mercury-bearing waste and must be disposed
of accordingly. Tartaric acid dust is a respiratory irritant.

%======================================================================
\begin{thebibliography}{99}

\bibitem{kinniburgh}
D.~G. Kinniburgh and W. Kosmus,
``Arsenic contamination in groundwater: some analytical considerations,''
\emph{Talanta}, vol.~58, pp.~165--180, 2002.
\url{https://nora.nerc.ac.uk/id/eprint/12300/1/ArsenicGround.pdf}

\bibitem{pat825}
Industrial Test Systems, Inc.,
``Arsenic analysis,'' US Patent 11{,}346{,}825~B2.
\url{https://image-ppubs.uspto.gov/dirsearch-public/print/downloadPdf/11346825}

\bibitem{pat300}
Industrial Test Systems, Inc.,
``Arsenic analysis'' (metal rate enhancers), US Patent 6{,}696{,}300.
\url{https://image-ppubs.uspto.gov/dirsearch-public/print/downloadPdf/6696300}

\bibitem{pat143}
``Method for the removal of sulfide interference in common tests for
arsenic,'' US Patent 6{,}576{,}143.
\url{https://image-ppubs.uspto.gov/dirsearch-public/print/downloadPdf/6576143}

\bibitem{rahman}
M.~M. Rahman, K. Fujinaga, Y. Seike and M. Okumura,
``A simple in situ visual and tristimulus colorimetric method for the
determination of trace arsenic in environmental water after its collection
on a mercury(II)-impregnated paper,''
\emph{Analytical Sciences}, vol.~20, no.~1, pp.~165--170, 2004.
\url{https://pubmed.ncbi.nlm.nih.gov/14753277/}

\bibitem{reddy}
R.~R. Reddy \emph{et al.},
``Evaluation of arsenic field test kits for drinking water: recommendations
for improvement and implications for arsenic affected regions such as
Bangladesh,''
\emph{Water Research}, vol.~170, art.~115325, 2020.
\url{https://www.sciencedirect.com/science/article/pii/S0043135419310991}

\bibitem{george}
C.~M. George \emph{et al.},
``Evaluation of an arsenic test kit for rapid well screening in
Bangladesh,''
\emph{Environmental Science \& Technology}, vol.~46, pp.~11213--11219, 2012.
\url{https://pubmed.ncbi.nlm.nih.gov/22866936/}

\bibitem{bird}
F.~C.~J. Bird, ``The Gutzeit test for arsenic,''
\emph{The Analyst}, vol.~24, p.~181, 1899.

\bibitem{km}
P. Kubelka and F. Munk,
``Ein Beitrag zur Optik der Farbanstriche,''
\emph{Zeitschrift f\"ur Technische Physik}, vol.~12, pp.~593--601, 1931.

\bibitem{crc}
\emph{CRC Handbook of Chemistry and Physics}, 84th ed., 2004 ---
dissociation constants for tartaric acid
($\mathrm{p}K_{a1} = 2.98$, $\mathrm{p}K_{a2} = 4.34$ at $25\,^\circ$C).

\bibitem{rsc}
``Review of analytical techniques for arsenic detection and determination in
drinking water,''
\emph{Environmental Science: Advances}, 2023, DOI 10.1039/D2VA00218C.

\bibitem{its}
Industrial Test Systems, Inc.,
Quick\texttrademark{} Arsenic Test Kit instruction booklet, part 481300-5.
\url{http://site.jjstech.com/pdf/ITS/481300-5\%20Instruction\%20Booklet.pdf}

\bibitem{nrt}
National Response Team,
``Quick Reference Guide: Arsine (SA),'' April 2017.
\url{https://nrt.org/sites/2/files/Arsine\%20SA\%20QRG\%20FINAL\%202017\%2004\%2025.pdf}

\bibitem{merckpaper}
Arsenic test paper (mercuric bromide paper), product insert ---
note that hydrogen sulfide also produces a yellow colour.
\url{https://www.ctlscientificsupply.com/prodinsert/90762insert.pdf}

\end{thebibliography}

\end{document}
